% !TITLE 近代
% !SUMMARY 旧诗新命
% !ORDER 0

\begin{intro}
近代中国，内忧外患交加，古老的诗歌传统面临着前所未有的挑战。

鸦片战争打开了中国的国门，也打破了士大夫“天朝上国”的迷梦。甲午战败、庚子事变、辛亥革命——每一次巨变都在诗歌中留下了深深的烙印。龚自珍的“我劝天公重抖擞，不拘一格降人才”，写的是一个先知先觉者对时代沉疴的焦虑；谭嗣同的“我自横刀向天笑，去留肝胆两昆仑”，写的是一个改革者面对死亡的从容；秋瑾的“一腔热血勤珍重，洒去犹能化碧涛”，写的是一个女革命家以身许国的决绝。

这一时期的诗歌，形式上仍然是旧体，精神上却已经是新的了。诗人们不再满足于个人的离愁别绪，而是把国家命运、民族危亡纳入了诗歌的视野。梁启超提出“诗界革命”，主张以旧风格写新意境；黄遵宪主张“我手写我口”，力求打破古典诗歌的程式化表达。这些尝试虽然还没有达到成熟，却为后来新诗的诞生埋下了伏笔。

王国维是一个特殊的例子。他的《人间词话》以西方美学观念重新审视中国古典诗词，提出了“境界说”，成为二十世纪中国文艺理论的重要起点。“昨夜西风凋碧树，独上高楼，望尽天涯路”——他用三句宋词概括了做学问的三个境界，这种打通中西、融会古今的尝试，正是近代知识分子精神的真实写照。

旧诗在近代没有死去，它只是换了一种方式继续活着：更沉重，更焦灼，更与时代同呼吸。
\end{intro}
